\begin{python0}
from solutions import *; clear()
\end{python0}
\sectionthree{Bubblesort and binary search}

\begin{ex}
Write a function \texttt{bubblesort()} that accepts an
array of integers and an integer for the size of the array and performs
bubblesort on the array so that the values are in ascending order. Test
your code: Create an array of 10 random integers, call the
\texttt{bubblesort()} function, and print the values of \texttt{x} in
\texttt{main()}. Test your function.
\end{ex}

\begin{ex}
Write a function
\texttt{\textbf{bubblesort\_from()}} that accepts an array \texttt{\textbf{x}} of integers and two integers \texttt{\textbf{start}} and \texttt{\textbf{end }}and performs bubblesort on from \texttt{\textbf{x[start]}} to \texttt{\textbf{x[end - 1]}} so that the values are in ascending order.
Test your code: Create an array of 10 random integers, call the \texttt{\textbf{bubble}\textbf{s}\textbf{ort\_}\textbf{from}\textbf{()}} function to sort the values using \texttt{\textbf{start=2}} and \texttt{\textbf{end=7}}, and print the values of \texttt{\textbf{x}} in \texttt{\textbf{main()}}. Test your function.
\end{ex}

\begin{ex}
Write a function \texttt{binarysearch()} that performs binary search on an array. Specifically, the function accepts an array \texttt{\textbf{x}} of integers, an integer \texttt{\textbf{x\_len}} for the size of the array, and an integer \texttt{\textbf{target}}. The function returns the index in \texttt{x} where \texttt{\textbf{target}} appears. The function assume that \texttt{\textbf{x}} is sorted in ascending order. Test your function.
\end{ex}

\begin{ex}
Write a function \texttt{binarysearch()} that accepts
an integer array \texttt{\textbf{x}}, integers \texttt{\textbf{start}} and
\texttt{\textbf{end}} and integer \texttt{\textbf{target}}, and then
performs binary search on an array \texttt{\textbf{x}} from index
\texttt{\textbf{start}} to index \texttt{\textbf{end - 1}} searching for
\texttt{\textbf{target. }}The function returns the index in \texttt{x} where
\texttt{\textbf{target}} appears. The function assume that
\texttt{\textbf{x}} is sorted in ascending order from index
\texttt{\textbf{start}} to index \texttt{\textbf{end - 1}}. Test your
function.
\end{ex}

